
Token log-probabilities are computed without additional cost during generation. Every autoregressive language model already computes them during the forward pass, and a substantial body of work [Kadavath et al., 2022; Fadeeva et al., 2024; Farquhar et al., 2024] treats the aggregated sequence-level statistic as a confidence proxy for hallucination detection. Yet the aggregation step itself --- the choice of which scalar to extract from the per-token log-probability sequence --- has not been the subject of systematic study. Most methods adopt one function (mean in CCP \citep{Fadeeva2024Fact}; sum within semantic clusters in Semantic Entropy \citep{Farquhar2024}) without testing whether alternative choices would perform better, and without a theoretical basis for expecting it to matter.

It matters by up to 12 AUROC points.

On TriviaQA (factual recall), taking the minimum token log-probability outperforms the mean by 5.6--11.9 AUROC points. On TruthfulQA (imitative falsehood), the mean outperforms the minimum by 11.1--12.0 AUROC points. The direction of advantage reverses cleanly between benchmark types, and this reversal is predicted by the structural nature of the hallucination being detected.

\textbf{The mechanism.} Recall-failure hallucinations --- where a model fails to retrieve a specific fact --- produce \textit{peaked} token probability distributions: the model generates function words with high probability but concentrates uncertainty at a single incorrect fact-token. The minimum log-probability captures this worst-case signal directly. Imitative-falsehood hallucinations --- where a model confidently generates an answer learned from training data that is factually wrong --- produce \textit{flat} high-probability distributions: no single token carries anomalous uncertainty, and the hallucination signal is distributed across all tokens. The mean log-probability integrates this diffuse signal that the minimum cannot resolve.

This mechanism is grounded in three independent layers of evidence. First, on TriviaQA, hallucinated responses show significantly higher token distribution peakedness than correct responses (peakedness ratio 2.936 vs. 2.533, two-sample t-test p = 0.002, n = 488), establishing Step 1 empirically for LLaMA-2-7B. Second, Spearman rank-correlation confirms the directional advantage without threshold assumptions: on all four (model $\times$ dataset) pairs, the direction of $\rho$(min) vs. $\rho$(mean) is consistent with the mechanism prediction, and all four bootstrap 95\% confidence intervals exclude zero. Third, AUROC threshold analysis quantifies the practical magnitude, with effect sizes 5.6--11.9 AUROC points on TriviaQA and 11.1--12.0 AUROC points on TruthfulQA.

An additional finding concerns the raw unnormalized log-probability sum. Contrary to the initial prediction that it would be the weakest aggregation, raw sum achieves the highest AUROC on TriviaQA for both models (0.895--0.896), exceeding both minimum (0.849--0.892) and mean (0.730--0.836), and exceeding the published Semantic Entropy AUROC ($\approx 0.79)$ \citep{Farquhar2024} at substantially lower inference cost. On TruthfulQA, raw sum is the weakest aggregation (0.417--0.459), producing a gap of approximately 0.44 AUROC between the two benchmark types for this statistic --- the largest single-aggregation contrast in the results.

\textbf{Contributions:}

\begin{enumerate}
\item \textbf{Mechanistic taxonomy.} Hallucination type (recall-failure vs. imitative-falsehood) is identified as the determinant of token probability distribution shape (peaked vs. flat), and distribution shape as the determinant of optimal aggregation function (min vs. mean). This is the first controlled ablation isolating aggregation function as the sole experimental variable on fixed factual QA splits, filling the gap identified by Liu et al. [2025].
\end{enumerate}

\begin{enumerate}
\item \textbf{Distributional evidence.} Direct experimental confirmation that recall-failure hallucinations produce significantly higher token distribution peakedness than correct responses at 7B model scale (p = 0.002, LLaMA-2-7B on TriviaQA), converting the mechanism's first step from assumption to measured fact.
\end{enumerate}

\begin{enumerate}
\item \textbf{Aggregation selection criterion.} A principled, label-free criterion for zero-cost hallucination detection: use minimum log-probability for factual recall tasks; use mean log-probability for imitative-falsehood tasks. Both require only a single forward pass with frozen model weights.
\end{enumerate}

\begin{enumerate}
\item \textbf{Unexpected finding.} Raw unnormalized log-probability achieves higher AUROC than both normalized aggregations on TriviaQA and surpasses published multi-sample baselines under this evaluation protocol, positioning sequence-level joint probability as a novel zero-cost detection signal warranting further investigation.
\end{enumerate}

The remainder of the paper is organized as follows. Section 2 positions this work within the hallucination detection literature. Section 3 describes the methodology, mechanism hypothesis, and four-hypothesis verification design. Section 4 presents the experimental setup. Section 5 reports results. Section 6 discusses implications and limitations. Section 7 concludes.

